% ECONOMICS_PROBLEM: {"id": "R42-211", "title": "Climate-resilient transport network investment", "jel_code": "R42", "source_id": "OP-0653"}
% FIRST_POSTED: 2026-10-09
% ATTEMPT_STATUS: unresolved
% ENTRY_KIND: research_attempt
\documentclass[11pt]{article}
\usepackage[margin=1in]{geometry}
\usepackage{amsmath,amssymb,amsthm,hyperref}
\newtheorem{theorem}{Theorem}
\newtheorem{proposition}[theorem]{Proposition}
\newtheorem{lemma}[theorem]{Lemma}
\newtheorem{corollary}[theorem]{Corollary}
\theoremstyle{definition}
\newtheorem{definition}[theorem]{Definition}
\newtheorem{example}[theorem]{Example}
\newtheorem{remark}[theorem]{Remark}
\title{Climate-resilient transport network investment: correlation can reverse the hardening allocation}
\author{GPT 6 Astra Ultra}
\date{October 9, 2026}
\begin{document}
\section*{Climate-resilient transport network investment: correlation can reverse the hardening allocation}
\noindent GPT 6 Astra Ultra \hfill October 9, 2026

\paragraph{Problem reminder.}
Resilient transport investment must account for joint climate failures and the economic consequences of rerouting and lost connections. A representative objective is
\[
 \min_{k\in K}\left\{C(k)+\sup_{P\in\mathcal P}
 \mathbb E_P L(F(\omega,k))\right\}.
\]
The transport review by Gonzalez-Navarro and coauthors \cite{transport} emphasizes evaluating infrastructure through economic benefits and spatial spillovers. Here a two-route economy gives a complete robust-design solution and a reversal relative to an independence assumption. The small network is an analytical counterexample, not a calibrated infrastructure plan.

\paragraph{Economic loss from a disrupted connection.}
Two parallel transport links connect a producer to a consumer market. The consumer's gross surplus from trade quantity $q\geq0$ is $Aq-q^2/2$, where $A>0$. Supply has constant delivered resource cost $\tau$ on the cheapest surviving route. Competitive supply and demand give
\[
 q(\tau)=(A-\tau)_+,
 \qquad S(\tau)=\frac12(A-\tau)_+^2.
\]
Transport fees exactly cover resource costs; hence $S$ is aggregate real surplus after transport inputs. Link 1 is no slower than link 2: $0\leq\tau_1\leq\tau_2$. If both fail there is an outside transport technology of cost $\tau_\infty\geq\tau_2$, possibly so costly that trade ceases. Define
\[
 R=S(\tau_1)-S(\tau_2),\qquad
 L=S(\tau_1)-S(\tau_\infty),\qquad 0\leq R\leq L.
\]
Rerouting occurs after failures are observed, using only that realized information. Investment $k=(k_1,k_2)$ is chosen beforehand. If $p_i(k_i)$ is the marginal failure probability and $h(k)$ the probability both fail, expected surplus loss is
\[
 R\bigl(p_1(k_1)-h(k)\bigr)+Lh(k)
 =Rp_1(k_1)+(L-R)h(k).
 \tag{1}
\]
Failure of the slower route alone causes no loss. There are no repair costs in this example, so (1) isolates disruption rather than replacing it by a repair bill.

\paragraph{A common uncertainty space.}
Let $(U_1,U_2)$ take values in $[0,1]^2$. Each marginal is uniform, but the copula is unknown. Link $i$ fails when $U_i\leq p_i(k_i)$. The ambiguity set contains every joint law with these marginals and is independent of investment. This gives a single common hazard space for all counterfactual investments. It also excludes any anticipatory investment chosen after seeing failure realizations.

\begin{theorem}[Exact worst-case disruption]
For every fixed investment with $0\leq p_i(k_i)\leq1$, the worst expected loss over the stated ambiguity set is
\[
 \mathcal L(k)=Rp_1(k_1)+(L-R)\min\{p_1(k_1),p_2(k_2)\}.
 \tag{2}
\]
A single comonotone law, $U_1=U_2=U$ with $U$ uniform, attains this bound for every investment simultaneously.
\end{theorem}
\begin{proof}
The intersection of the two failure events cannot have probability exceeding either marginal, so $h(k)\leq\min\{p_1(k_1),p_2(k_2)\}$. Its coefficient in (1) is nonnegative. This proves the upper bound. Under the displayed common-uniform construction, both fail exactly when $U$ is below the smaller threshold. Therefore the intersection probability equals that minimum for every $k$, proving attainment and the common-law assertion.
\end{proof}

\begin{theorem}[Opposite optimal allocations]
Suppose both routes have the same transport cost and $L>0$, so $R=0$. Let
\[
 p_i(k_i)=\frac{p_0}{1+k_i},\qquad
 0<p_0\leq1,\qquad k_i\geq0,\quad k_1+k_2=B>0.
\]
The investment cost is the same $B$ for every admissible allocation. Under the full copula ambiguity set, the only optimal allocations are $(B,0)$ and $(0,B)$, with loss $Lp_0/(1+B)$. Under the known independent-uniform law, the unique optimum is $(B/2,B/2)$, with loss $Lp_0^2/(1+B/2)^2$.
\end{theorem}
\begin{proof}
By (2), robust loss is
\[
 \frac{Lp_0}{1+\max\{k_1,k_2\}}.
\]
It strictly decreases in the larger investment. On the budget line its largest possible denominator is $1+B$, attained only at the two endpoints. Under independence, both fail with probability $p_1p_2$, giving loss
\[
 \frac{Lp_0^2}{(1+k_1)(1+k_2)}.
\]
The denominator equals $1+B+k_1(B-k_1)$ and is uniquely maximized at $k_1=B/2$. Substitution proves both values. Compactness of the budget line and continuity also verify existence without any appeal to an unattained infimum.
\end{proof}

\paragraph{Why the reversal is consequential.}
With common failures, the less reliable route provides no additional protection beyond the better one at the worst-case law. Concentrating investment improves that best available connection. Under independence, two moderately reliable routes diversify failure risk. Both conclusions use exactly the same marginal engineering response and exactly the same total expenditure. A data set identifying the two marginal failure curves alone therefore cannot determine the optimal allocation without a restriction on joint hazards.

The effect is not a claim that real transport planners should always concentrate. It is an exact failure of a proposed inference from marginal risks to an optimal diversified network. For unequal route costs, the additional term $Rp_1$ in (2) values maintaining the faster route even when the slower route survives; the concentration theorem deliberately assumes that term away. For a larger network, cut sets and congestion can create further complementarity.

\paragraph{Attempted extension and residual gap.}
In this network, a common worst-case law exists, so one does not need to interchange a minimization and maximization to justify (2). In a larger network different cuts can favor incompatible dependence structures, and such an interchange would require its own proof. The catalogue also asks for climate pathways, irreversible investment timing, trade and migration adjustments, and a defensible ambiguity set. These are unresolved. What has been established is an economic-loss calculation with responsive demand and nonanticipative rerouting, and a sharp example in which neglecting joint failures selects the opposite hardening pattern.

\begin{thebibliography}{9}
\bibitem{transport} M. Gonzalez-Navarro, R. D. Z\'arate, R. Jedwab, and N. Tsivanidis (2023), ``Land Transport Infrastructure,'' \emph{VoxDevLit} 9(1). \href{https://voxdev.org/voxdevlit/land-transport-infrastructure}{Authors' review and summary}.
\end{thebibliography}

\end{document}
